\section{Тепловые процессы в среде}

Перед изучением термоупругости необходимо ввести понятия теплопереноса и температурных деформаций. Эти эффекты будут затем связаны с механикой сплошной среды.

\subsection{Базовая теплопроводность}

В изотропной среде плотность теплового потока определяется законом Фурье:
\begin{equation}
\mathbf{q} = -k \nabla T ,
\end{equation}
где $k$ --- коэффициент теплопроводности, $T$ --- температура.  

Закон Фурье вместе с законом сохранения энергии приводит к тепловому (диффузионному) уравнению:
\begin{equation}
\rho c_p\, \frac{\partial T}{\partial t} = k \nabla^2 T,
\label{eq:heat}
\end{equation}
где $c_p$ --- теплоёмкость при постоянном давлении.  

Решения~(\ref{eq:heat}) описывают температурные поля, их распространение и затухание (тепловая диффузия) \cite{OpenStax_HE, MIT_heat}.


\subsection{Температурные деформации}

При нагревании материал стремится изменять длину. Для малых температурных приращений линейная деформация пропорциональна температуре:
\begin{equation}
\varepsilon_{th} = \alpha\,\Delta T ,
\end{equation}
где $\alpha$ --- коэффициент линейного температурного расширения.  

В трёхмерном виде температурная деформация записывается тензорно:
\begin{equation}
\varepsilon^{th}_{ij}= \alpha \Delta T\,\delta_{ij},
\end{equation}
где $\delta_{ij}$ --- символ Кронекера.

Температурные расширения приводят к внутренним упругим напряжениям, если расширение материала ограничено внешними условиями (например, в твёрдом теле с закреплёнными границами). Это ключевой механизм будущей термоупругой связи \cite{Feynman_vol2, Purdue_MM}.


\section*{Цель раздела}

Понять:
\begin{itemize}
\item как распространяется тепло,
\item как температура формирует температурные поля,
\item как температура вызывает упругие деформации.
\end{itemize}
Эти эффекты затем объединяются с уравнениями Ламе, формируя термоупругие уравнения (тепло + упругость).


\section*{Список источников}

\begin{thebibliography}{9}

\bibitem{OpenStax_HE}
OpenStax. \textit{Heat and Heat Transfer}. Physics.  
URL: \url{https://openstax.org/books/physics/pages/14-introduction}

\bibitem{MIT_heat}
MIT OpenCourseWare. \textit{Heat Transfer and Thermal Diffusion}. Lecture notes.  
URL: \url{https://ocw.mit.edu/courses/2-51-intermediate-heat-and-mass-transfer-fall-2008/lecture-notes/}

\bibitem{Feynman_vol2}
Feynman, R. \textit{The Feynman Lectures on Physics, Vol.2}. Thermal physics section.  
Open online edition: \url{https://www.feynmanlectures.caltech.edu/II_toc.html}

\bibitem{Purdue_MM}
Purdue University. \textit{Thermal Effects in Materials (lecture notes)}.  
URL: \url{https://engineering.purdue.edu/~materials/thermal_effects.pdf}

\end{thebibliography}
